\documentclass[12pt,a4paper]{article}

% --- Paquetes de diseño y fuentes ---
\usepackage[utf8]{inputenc}
\usepackage[spanish, es-tabla]{babel}
\usepackage{amsmath, amsfonts, amssymb, amsthm}
\usepackage{graphicx}
\usepackage{geometry}
\usepackage{hyperref}
\usepackage{booktabs}
\usepackage{fancyhdr}
\usepackage{setspace}

% Configuración de márgenes
\geometry{top=2.5cm, bottom=2.5cm, left=2.5cm, right=2.5cm}
\onehalfspacing

% Cabecera y Pie de página
\pagestyle{fancy}
\fancyhf{}
\rhead{\small Semillero de Aprendizaje por Refuerzo}
\lhead{\small IA para Trading y Sentimiento}
\cfoot{\thepage}
\renewcommand{\headrulewidth}{0.4pt}

% Título y Autores
\title{
    \vspace{-1cm}
    \textbf{\Large IA y Emociones: Desarrollo de un Agente de Aprendizaje por Refuerzo para el Mercado de Criptomonedas} \\
    \vspace{0.5cm}
    \large Universidad Externado de Colombia \\
    \text{Ciencia de Datos -- Departamento de Matemáticas}
}
\author{
    \textbf{Julian Duarte} \quad \textbf{Julian Jimenez} \\
    \medskip
    \small \textit{Tutores:} \\
    \small Prof. Julián Fernando Sánchez \quad Prof. Juan Carlos Urueña \quad Prof. Camilo Castillo
}
\date{Marzo de 2026}

\begin{document}

\maketitle

\begin{abstract}
    Este trabajo presenta el desarrollo de un Agente de Aprendizaje por Refuerzo (IA) diseñado para operar en el mercado de criptomonedas. A diferencia de los modelos tradicionales, nuestra propuesta integra el "estado de ánimo" colectivo mediante un índice de sentimiento, permitiendo al agente tomar decisiones más informadas. El objetivo es maximizar la rentabilidad acumulada gestionando el riesgo de forma inteligente. Explicamos paso a paso cómo la IA observa el mercado, elige sus acciones y aprende de sus aciertos y errores.
\end{abstract}

\section{Introducción y Objetivo del Agente}
El mercado de criptomonedas se caracteriza por su alta volatilidad e influencia psicológica. El objetivo principal de nuestro agente es encontrar una política de trading óptima que maximice su **capital final acumulado $(\sum \gamma^t R_t)$**. No buscamos solo ganar dinero rápido, sino aprender a proteger el capital en momentos de incertidumbre y aprovechar las oportunidades de crecimiento, equilibrando el beneficio con la estabilidad.

\section{Configuración de Datos y Entorno}
Para que nuestra IA aprenda en un entorno realista, hemos definido las siguientes especificaciones técnicas:
\begin{itemize}
    \item \textbf{Activo de Interés:} BNB/USDT (Binance Coin operada contra Tether).
    \item \textbf{Fuente de Datos:} API de Binance para precios (OHLCV) y \textit{alternative.me} para el sentimiento.
    \item \textbf{Temporalidad:} Velas de 1 hora (1h). Esta frecuencia permite capturar movimientos intradía sin el ruido excesivo de los minutos.
    \item \textbf{Periodo de Análisis:} Enero de 2023 a Diciembre de 2024. Este intervalo se eligió por su **relevancia estratégica**, ya que abarca tanto la recuperación post-crisis de 2023 como el ciclo alcista de 2024, ofreciendo una variedad de regímenes de mercado para el entrenamiento del agente.
\end{itemize}

\section{Sentimiento de Mercado: El Pulso Emocional}
Hemos seleccionado el \textbf{Crypto Fear \& Greed Index} como nuestra fuente única de sentimiento. Este índice traduce múltiples indicadores sociales y de mercado en un número de 0 a 100:
\begin{itemize}
    \item \textbf{0 - 24 (Miedo Extremo):} Posibles oportunidades de compra (mercado infravalorado).
    \item \textbf{25 - 49 (Miedo):} Cautela general.
    \item \textbf{50 - 74 (Codicia):} Optimismo creciente.
    \item \textbf{75 - 100 (Codicia Extrema):} Peligro de corrección (mercado sobrevalorado).
\end{itemize}
El agente recibe este valor diariamente y lo alinea con los precios horarios para refinar su toma de decisiones.

\section{El Estado: ¿Qué observa la IA? ($S_t$)}
El "Estado" es la brújula del agente. En cada hora, la IA observa un vector $S_t$ compuesto por:
\begin{enumerate}
    \item \textbf{Régimen de Mercado ($M_t$):} ¿El precio está subiendo o bajando respecto a la hora anterior?
    \item \textbf{Sentimiento ($Sent_t$):} El nivel actual del índice de Miedo y Codicia.
    \item \textbf{Indicador de 'Cansancio' (RSI):} El Índice de Fuerza Relativa nos dice si el activo está sobrecomprado o sobrevendido.
    \item \textbf{Inventario ($I_t$):} Una variable binaria (0 si tenemos efectivo, 1 si tenemos el activo).
\end{enumerate}

\section{Acciones y Recompensas (Premio y Castigo)}
El agente puede elegir entre \textbf{Comprar}, \textbf{Vender} o \textbf{Mantener}. Para aprender, recibe una recompensa $R_{t+1}$ calculada bajo tres pilares:
\begin{equation}
    R_{t+1} = \underbrace{I_{t+1}r_{t+1}}_{\text{Ganancia}} - \underbrace{c|I_{t+1} - I_t|}_{\text{Comisiones}} - \underbrace{\lambda \sigma_t}_{\text{Penalización por Riesgo}}
\end{equation}
Donde $c$ representa los costos de transacción y $\lambda \sigma_t$ es un castigo proporcional a la volatilidad ($\sigma$). Esto enseña al agente que las ganancias con alto riesgo son menos valiosas que las ganancias estables.

\section{Cerebro de la IA: Q-Learning Tabular}
Para resolver este problema, utilizamos **Q-Learning**, un algoritmo que crea una "Tabla de Sabiduría" (Q-Table). En esta tabla, el agente anota qué tan buena es cada acción para cada situación posible.
\begin{itemize}
    \item \textbf{Aprendizaje:} Si una acción lleva a una buena recompensa, su valor en la tabla aumenta.
    \item \textbf{Decisión:} Ante una situación nueva, el agente consulta su tabla y elige la acción con el valor más alto acumulado históricamente.
\end{itemize}
Este enfoque es ideal para el póster porque permite explicar visualmente cómo la experiencia se convierte en decisiones matemáticas.

\bibliographystyle{plain}
\begin{thebibliography}{9}
    \bibitem{sutton} Sutton, R. S., \& Barto, A. G. (2018). \textit{Reinforcement learning: An introduction}. MIT press.
\end{thebibliography}

\end{document}
